% Generated draft: requires mathematical and source review.
\documentclass[11pt]{article}
\usepackage[margin=25mm]{geometry}
\usepackage{fontspec}
\setmainfont{Latin Modern Roman}
\usepackage{amsmath,amssymb,mathtools,unicode-math}
\setmathfont{Latin Modern Math}
\usepackage{enumitem,needspace,xurl,hyperref}
\setlength{\parindent}{0pt}
\setlength{\parskip}{5pt}
\setlength{\emergencystretch}{3em}
\newcommand{\N}{\mathbb{N}}
\newcommand{\Z}{\mathbb{Z}}
\newcommand{\Q}{\mathbb{Q}}
\newcommand{\R}{\mathbb{R}}
\newcommand{\C}{\mathbb{C}}
\newcommand{\F}{\mathbb{F}}
\newcommand{\A}{\mathbb{A}}
\newcommand{\PP}{\mathbb{P}}
\newcommand{\E}{\mathbb{E}}
\newcommand{\eps}{\varepsilon}
\newcommand{\dd}{\,\mathrm{d}}
\begin{document}
\section*{UnsolvedMath ID: 11100045}
\textbf{Source problem code:} catalogue entry
\subsection{Definitions and mathematical statement}
\textbf{Equivariant homotopy problem 3 (Hovey).} Develop equivariant stable homotopy theory without the usual restrictions on the acting group, explicitly including infinite discrete groups. Determine the required changes to the existing setup rather than assume that a compact-group construction applies unchanged.

A precise version should specify, for the intended groups $G$, categories of equivariant spectra and stable homotopy categories $\mathrm{Ho}(\mathrm{Sp}^G)$, their weak equivalences, and which equivariant operations and comparison properties are to be retained. The excerpt gives no formal axioms or a fixed model. In particular, it does not say whether ``all groups'' means abstract groups viewed discretely or a stated class of topological groups, nor whether all actions or only proper actions with a prescribed isotropy family are in scope.

For any discrete $G$, spectra with an ordinary $G$-action give a naive equivariant category, but this alone does not specify the genuine fixed-point, transfer, representation-grading, or change-of-group structure sought in an extension of usual equivariant stable homotopy theory. Conversely, transfers available in a compact or finite-index setting must not be required for arbitrary infinite-index subgroup inclusions without a suitable replacement. Deciding which structure survives is part of the requested reformulation.

Hovey credits Mike Mandell with ideas and emphasizes infinite discrete groups as the setting of important problems such as the Novikov conjecture. This is an application motivation, not a request that a proposed framework already prove Novikov. The summary mentions global and spectral-Mackey frameworks and assembly maps, but supplies no precise construction, group scope, or theorem identifying them with a complete answer to the source's unrestricted program.
\subsection{Short English statement}
Extend equivariant stable homotopy theory to unrestricted acting groups, especially infinite discrete groups, while specifying which genuine equivariant structures and operations are available.
\subsection{Sources}
\begin{itemize}
\item Local UnsolvedMath registry, permanent ID 11100045, source code catalogue entry; attributed to Mark Hovey, no proposal year supplied.
\item \emph{Hovey's Algebraic Topology Problem List}, Equivariant homotopy, item 3; surviving mirror \url{https://www-users.cse.umn.edu/~tlawson/hovey/}. The registry records HTML-ordered-list extraction and access on 2026-07-29; the mirror has 65 top-level items on nine category pages.
\item Mike Mandell is named as having relevant ideas, without a paper or theorem. The Novikov conjecture is cited as motivation. The summary's global/spectral-Mackey and assembly-map references are unnamed and bibliographically incomplete.
\end{itemize}
\end{document}
